%
% 44-hauptwert.tex -- Hauptwert
%
% (c) 2026 Prof Dr Andreas Müller
%
\subsection{Hauptwert
\label{buch:integration:residuen:subsection:hauptwert}}
Die Berechnung des Integrals in
Abschnitt~\ref{buch:integration:residuen:section:beispiel}
hat einen Weg verwendet, der den Nullpunkt in einem grossen
Bogen umrundet hat.
So wurde das Integral zwischen $-R$ und $R$ zu einem Wegintegral
über eine geschlossene Kurve vervollständigt, wodurch sich der
Integralsatz anwenden liess.
Auf diese Weise konnte ein uneigentliches Integral von 
$1/(x^2+1)$ über die ganze reelle Achse berechnet werden.

Auch das Integral
\input{chapters/030-integration/fig/fig-sinxx.tex}%
\begin{equation}
I
=
\int_{-\infty}^\infty \frac{\sin x}{x}\,dx
\label{buch:integration:residuen:eqn:sinxxintegral}
\end{equation}
ist uneigentlich.
Da der Integrand mit Nullstellen bei ganzzahligen Vielfachen
von $\pi$ oszilliert, nimmt die Stammfunktion
\[
x
\mapsto
F(x)
=
\int_{0}^{x} \frac{\sin t}{t}\,dt
\]
lokale Extrema bei diesen Nullstellen an
(Abbildung~\ref{buch:integration:fig:sinxx}).
Jeder Wert liegt daher zwischen Werten
\[
F\biggl(
\biggl\lfloor\frac{x}{\pi}\biggr\rfloor\pi
\biggr)
=
\int_{0}^{\lfloor\frac{x}{\pi}\rfloor\pi}
\frac{\sin t}{t}\,dt
\qquad
\text{und}
\qquad
F\biggl(
\biggl(\biggl\lfloor\frac{x}{\pi}\biggr\rfloor+1\biggr)\pi
\biggr)
=
\int_{0}^{\left(\lfloor\frac{x}{\pi}\rfloor+1\right)\pi}
\frac{\sin t}{t}\,dt
\]
der Stammfunktion an den nächstliegenden ganzzahligen Vielfachen von $\pi$.
Die Konvergenz des uneigentlichen Integrals ist damit gesichert.
Trotzdem bleibt die Berechnung des Integralwertes offen, da es
keine geschlossene Form für die Stammfunktion gibt.

Um den Werte des Integrals $I$ mithilfe eines Wegintegrals zu berechnen,
muss ein komplexer Integrand und ein passender Weg gefunden werden.
Dabei ist zu berücksichtigen, dass der Integrand an der Stelle $x=0$
eine Singularität hat.
\input{chapters/030-integration/fig/fig-hauptwert.tex}%
Als Integrand bietet sich
\[
f(z)
=
\frac{e^{iz}}{z}
\]
an.
Während der ursprüngliche reelle Integrand an der Stelle $x=0$ stetig
ist, hat $f(z)$ dort einen Pol.
Als Weg kann man der in
Abbildung~\ref{buch:integration:residuen:fig:hauptwert}
dargestellt Weg verwendet werden.
Er setzt sich zusammen aus den beiden Segmenten $[\varepsilon,R]$
und $[-R,-\varepsilon]$ und den Kreisbögen mit Radien $R$ und $\varepsilon$,
die die Segmente zu einem geschlossenen Weg verbinden.
Im Inneren dieses Weges ist $f(z)$ eine holomorphe Funktion.
Das gesuchte Integral ist also der Teil über die beiden Segmente
$\gamma_\pm$ im Grenzwert $R\to\infty$ und $\varepsilon\to 0$:
\begin{align*}
I
&=
\lim_{R\to\infty}
\lim_{\varepsilon\to 0}
\biggl(
\int_{\gamma_+}
f(z)\,dz
+
\int_{\gamma_-}
f(z)\,dz
\biggr)
\intertext{Aus dem Integralsatz kann man dann schliessen, dass
das Integral auch}
&=
-
\lim_{R\to\infty}
\int_{\gamma_R} f(z)\,dz
-
\lim_{\varepsilon\to 0}
\int_{\gamma_{\varepsilon}} f(z)\,dz
\end{align*}
ist.
Bei dieser Berechnung berechnet man das Integral entlang der
reellen Achse aber mit symmetrischen Segmenten.
Die Definition eines uneigentlichen Integrals verlangt, dass für
die Grenzen unabhängig voneinander ein Grenzübergang durchgeführt
wird.
Der symmetrische Grenzwert heisst auch der \emph{Hauptwert} des
Integrals und wird im Kapitel~\ref{chapter:hauptwert} ausführlicher 
dargestellt.




